\documentclass[10pt]{article}   % document properties
\usepackage[spanish]{babel}     % spanish package
\usepackage[utf8]{inputenc}     % codificacion de caracteres utf-8
\usepackage[top=2.3cm,bottom=3.5cm,right=2.5cm,left=2.5cm]{geometry}   % modify margins 
\usepackage{multicol}           % multiple columsn
\usepackage{graphicx}           % manage images 
\usepackage{url}                % manage aesthetic links
\usepackage{hyperref}           % internal links and more
\usepackage{array}              % customize tables         
%\usepackage{tcolorbox}          % colorful boxes
\usepackage{caption}            % improve figures & tables captions
\usepackage{tabularx}           % table advanced (sized)
\usepackage{lipsum}             % lipsum text
\usepackage{fancyhdr}           % headers and footers
\usepackage{xcolor, colortbl}   % define customize colors & table colore
\usepackage{listings}           % show & customize code
%%%%%%%%%%%%%%%%%%%%%%%%%%%%%%  extra definitions  %%%%%%%%%%%%%%%%%%%%%%%%%%%%%%%
\lstdefinestyle{ascii-tree}{
	literate={├}{|}1 {─}{--}1 {└}{+}1 
}
\newcolumntype{M}[1]{>{\centering\arraybackslash}m{#1}}     % cell configuration
\pagestyle{fancy}               % choosing fancy style
\fancyhf{}                      % delete any footer or header
\renewcommand{\headrulewidth}{0pt} % delete header line
% para el codigo fuente
\definecolor{dkgreen}{rgb}{0,0.6,0}
\definecolor{gray}{rgb}{0.5,0.5,0.5}
\definecolor{graya}{rgb}{0.97,0.97,0.97}
\definecolor{mauve}{rgb}{0.58,0,0.82}
\definecolor{codebackground}{rgb}{0.95, 0.95, 0.92}
\definecolor{tablebackground}{rgb}{0.8, 0, 0}
\lstset{
      frame=tb,
      language=bash,
      aboveskip=3mm,
      belowskip=3mm,
      showstringspaces=false,
      columns=flexible,
      basicstyle={\small\ttfamily},
      numbers=left,
      numberstyle=\tiny\color{gray},
      keywordstyle=\color{blue},
      commentstyle=\color{dkgreen},
      stringstyle=\color{mauve},
      breaklines=true,
      breakatwhitespace=true,
      tabsize=3,
      backgroundcolor= \color{codebackground}
}
%%%%%%%%%%%%%%%%%%%%%%%%%%%%%%%%%%%  variables  %%%%%%%%%%%%%%%%%%%%%%%%%%%%%%%
\newcommand{\itemCourse}{Nombre del Curso del Informe}
\newcommand{\itemTheme}{Titulo de tu Informe}
\newcommand{\itemPracticeNumber}{0?}
\newcommand{\itemAcademic}{2023 - B}
\newcommand{\itemSemester}{VIII} %Romanos de preferencia
\newcommand{\itemDate}{?? / ?? / 2023}
\newcommand{\itemHour}{23:59 PM}

\newcommand{\itemStudentA}{Participante 1}
\newcommand{\itemStudentB}{Participante 2}
\newcommand{\itemStudentC}{Participante 3}
\newcommand{\itemStudentD}{Participante 4}

\newcommand{\itemTeacher}{Dr(a). Docente a cargo}
\newcommand{\itemUniversity}{Universidad Nacional de San Agustín de Arequipa}
\newcommand{\itemFaculty}{Facultad de Ingeniería de Producción y Servicios}
\newcommand{\itemDepartment}{Departamento Académico de Ingeniería de Sistemas e Informática}
\newcommand{\itemSchool}{Escuela Profesional de Ingeniería de Sistemas}
%%%%%%%%%%%%%%%%%%%%%%%%%%%%%%%%%%%   header  %%%%%%%%%%%%%%%%%%%%%%%%%%%%%%%%%%

\fancyhead[C]{
	\begin{tabularx}{\textwidth}{|M{3.8cm}|M{7.67cm}|M{3.8cm}|}
		\hline
		\includegraphics[scale=0.45]{img/logo_episunsa.png} &\cellcolor{graya}{\textbf{\footnotesize\itemUniversity}\par\textbf{\footnotesize\itemFaculty}\par\textbf{\footnotesize\itemSchool}} & \vspace{0.16cm}\includegraphics[scale=0.52]{img/logo_abet.png} \\
		\hline
		\multicolumn{3}{|c|}{\footnotesize \textbf{Formato:} Guía de Práctica de Laboratorio / Talleres / Centros de Simulación}\\
		\hline
		\cellcolor{graya}{\textbf{\footnotesize Aprobación: 2022/03/01}} & \textbf{\footnotesize Código: GUIA-PRLD-001} & \cellcolor{graya}{\textbf{\footnotesize Página: \thepage}} \\
		\hline
	\end{tabularx}}

\setlength{\headheight}{66pt}   % Ajusta la altura del encabezado
\setlength{\headsep}{16pt}      % Ajusta la separación entre el encabezado y el texto
\setlength{\footskip}{0pt}      % Ajusta la separación entre el final del texto y el pie de página

\begin{document}
    \vspace*{0cm}	
    \begin{center}	
        \fontsize{17}{17} \Large{\textbf{INFORME DE LABORATORIO}}
    \end{center}
    
    \begin{table}[h!]
        \renewcommand{\arraystretch}{1.7}
        \footnotesize
        \begin{tabular}{|m{2.4cm}|m{2.1cm}|m{2.4cm}|m{2cm}|m{2.64cm}|m{2.42cm}|}\hline 
            \rowcolor{tablebackground}
            \multicolumn{6}{|c|}{\textbf{\large\color{white} INFORMACION BASICA}}\\ \hline
            {\cellcolor{graya}{ASIGNATURA:}} & \multicolumn{5}{l|}{\itemCourse}\\ \hline 
            \cellcolor{graya}{TITULO DE LA PRACTICA:} & \multicolumn{5}{l|}{\itemTheme}\\ \hline 
            \cellcolor{graya}{NUMERO DE LA PRACTICA:} & \itemPracticeNumber & \cellcolor{graya}{AÑO LECTIVO:} & \itemAcademic & \cellcolor{graya}{N° SEMESTRE:} & \itemSemester\\ \hline 
            \cellcolor{graya}{FECHA DE \par PRESENTACION:} & \itemDate & \cellcolor{graya}{HORA DE \par PRESENTACION:} & \multicolumn{3}{l|}{\itemHour} \\ \hline 
            \multicolumn{4}{|l|}{\begin{minipage}{8cm}
                \vspace{0.5em}
                INTEGRANTE (s):
                \begin{itemize}
                    \setlength{\itemsep}{0pt}
                    \setlength{\parskip}{0pt}
                    \setlength{\parsep}{0pt}
                    \item \itemStudentA
                    \item \itemStudentB
                    \item \itemStudentC
                    \item \itemStudentD
                \end{itemize}
                \vspace{0em} % Espaciado ajustable según necesidad
            \end{minipage}} & \cellcolor{graya}{NOTA:} & \\ \hline 
            \multicolumn{6}{|l|}{\begin{minipage}{8cm}
                \vspace{0.5em} % Espaciado ajustable según necesidad
                DOCENTE (s):
                \begin{itemize}
                    \setlength{\itemsep}{0pt}
                    \setlength{\parskip}{0pt}
                    \setlength{\parsep}{0pt}
                    \item \itemTeacher
                \end{itemize}
                \vspace{0em} % Espaciado ajustable según necesidad
            \end{minipage}}\\ \hline 	
        \end{tabular}
    \end{table}
    \normalsize
    
    \section{Introducción}

\%El primer párrafo presenta información general sobre el tema en estudio o problema que se trata de resolver. Debe resaltar la importancia del tema tratado. y su impacto\%

\%El Segundo párrafo resalta las investigaciones previas\%

\%El tercer párrafo resalta el trabajo desarrollado , que se ha hecho (mencionar los casos) y como se hizo\%

\%El cuarto párrafo presenta la estructura del artículo. Siempre inicia con las palabras: ``El resto del artículo está organizado de la siguiente manera...'' \%


\section{Tema 1 – Realidad Aumentada}

\subsection{Definición y conceptos fundamentales}
Use automatic hyphenation and check spelling and grammar. Use high resolution (300dpi or above) figures, plots, drawings and photos for best printing result.

\subsection{Enfoques terapéuticos}
Use automatic hyphenation and check spelling and grammar. Use high resolution (300dpi or above) figures, plots, drawings and photos for best printing result.

\subsection{Desafíos y limitaciones}
Use automatic hyphenation and check spelling and grammar. Use high resolution (300dpi or above) figures, plots, drawings and photos for best printing result.


\section{Tema 2 – Marco Teórico}
\subsection{Subtema 1}
Use automatic hyphenation and check spelling and grammar. Use high resolution (300dpi or above) figures, plots, drawings and photos for best printing result.

\subsection{Subtema 2}
Use automatic hyphenation and check spelling and grammar. Use high resolution (300dpi or above) figures, plots, drawings and photos for best printing result.

\subsection{Subtema 3}
Use automatic hyphenation and check spelling and grammar. Use high resolution (300dpi or above) figures, plots, drawings and photos for best printing result.

\section{Metodología}

\subsection{Caso 1}

En el trabajo de Coronado et al. [10], se presenta el uso de la tecnología de Realidad Virtual (RV) en un sistema de control de balance para pacientes con Parkinson. El sistema utiliza el dispositivo Oculus, el cual proporciona una experiencia inmersiva a través de una simulación de bicicleta en un entorno virtual con un camino, árboles y obstáculos que deben ser evitados mediante el movimiento del cuerpo hacia la izquierda o derecha, respectivamente (Figura 4). Para la captura de los movimientos corporales, se integró el controlador Leap Motion Controller. El estudio se llevó a cabo en el Hospital Honorio Delgado, específicamente en el área de neurología.\\

Diez pacientes con enfermedad de Parkinson en estados iniciales participaron en el estudio, compuesto por cinco hombres y cinco mujeres, con edades comprendidas entre 55 y 78 años. El protocolo del estudio consistió en diez sesiones de intervención, realizadas tres veces por semana, con una duración de 30 minutos cada sesión.\\

Antes y después de las sesiones, se evaluó el desempeño de los pacientes utilizando pruebas clínicas, específicamente el Test de la Batería de Evaluación del Movimiento (BOT-2) [11] y el Test de Bloques y Cajas (Block and Box Test) [12]. Además, se recopiló la opinión de los pacientes sobre la usabilidad del sistema mediante cuestionarios de evaluación de satisfacción del usuario (User Satisfaction Evaluation Questionnaire, USEQ) [13], cuestionarios de experiencia del usuario (User Experience Questionnaire, UEQ) [14] y el Modelo de Aceptación Tecnológica (Technological Acceptance Model, TAM) [15].\\

Los resultados obtenidos mostraron una mejora significativa en la subprueba del balance del Test BOT-2 en ocho de los participantes, evidenciando los beneficios del sistema de control de balance basado en Realidad Virtual. Dos participantes no mostraron mejoría en esta área específica. En cuanto a las opiniones de los pacientes, cinco de ellos manifestaron que disfrutaron de la experiencia proporcionada por el sistema RV, mientras que los otros cinco experimentaron mareos después de utilizar el dispositivo. Los resultados del cuestionario UEQ revelaron una valoración positiva por encima de los 4 puntos, lo cual indica una alta aceptación del sistema por parte de los usuarios.


\subsection{Caso 2}
Se describe el contexto (institución y lugar), tiempo que duró la experiencia, cantidad de participantes, resultados obtenidos, software/hardware/metodología utilizados (según corresponda)

\subsection{Caso 3}
Se describe el contexto (institución y lugar), tiempo que duró la experiencia, cantidad de participantes, resultados obtenidos, software/hardware/metodología utilizados (según corresponda)

\subsection{Caso 4}
Se describe el contexto (institución y lugar), tiempo que duró la experiencia, cantidad de participantes, resultados obtenidos, software/hardware/metodología utilizados (según corresponda)
 
\section{Resultados y Discusión}
En el análisis comparativo de los cuatro casos de uso de realidad aumentada para rehabilitación en motor fino, se identificaron semejanzas y diferencias relevantes.

\subsection{Semejanzas}
Use automatic hyphenation and check spelling and grammar. Use high resolution (300dpi or above) figures, plots, drawings and photos for best printing result.

\subsection{Diferencias}
Use automatic hyphenation and check spelling and grammar. Use high resolution (300dpi or above) figures, plots, drawings and photos for best printing result.

\subsection{Discusión}
Use automatic hyphenation and check spelling and grammar. Use high resolution (300dpi or above) figures, plots, drawings and photos for best printing result.

\subsection{Propuesta}
Use automatic hyphenation and check spelling and grammar. Use high resolution (300dpi or above) figures, plots, drawings and photos for best printing result.

\section{Conclusiones} % Desarrolla tu contenido dentro de desarrollo
    %\input{conclusiones} % agregar las partes necesarias
    
\end{document}